\subsection*{7. Quality, Perception, and Economics Are Separate Constraints}

A company may produce a genuinely superior product but fail because its target customers cannot recognize the improvement. It may produce a convincing sales story but fail to deliver. Or it may deliver exceptionally well while destroying its own economics through unconstrained service labor.

For each proposed difference, test 3 questions:

\begin{enumerate}
    \item \textbf{Buyer relevance:} Compared with the best real alternative, does this change something the customer values?
    \item \textbf{Credibility and delivery:} Can a skeptical customer inspect reasons to believe it, and can we fulfill it consistently?
    \item \textbf{Economics and resilience:} Can we create and capture value after acquiring, serving, supporting, and retaining the customer? What happens when rivals respond?
\end{enumerate}

Differentiation matters when it causes the right customer to choose you, is demonstrably real after purchase, and can be economically repeated. Durable protection is a separate, longer-run question.

\subsection*{8. Why Pure Superiority Is Not Always the Winning Mechanism}

A common error in business theory is to assume the best-performing product wins. For low-involvement consumer goods, buyers often prefer the accessible, familiar, sufficiently good option. They do not run a scientific tournament between all biscuits or household cleaners. Recognition, availability, pack size, habit, and retail economics can be decisive. Research associated with the Ehrenberg-Bass Institute challenges the universal necessity of perceived differentiation for brand growth [R3].

For a hospital purchasing complex equipment, by contrast, documentation, safety, total cost of ownership, regulatory compliance, and maintenance capacity dominate. A billboard campaign cannot compensate for a failure in technical approval.

This suggests a governing diagnostic: \emph{How is a decision actually made in this category?} Determine whether it is habitual or deliberative, frequent or rare, low-risk or high-stakes, individual or committee-based, observable or difficult to verify, and easy or costly to reverse. That predicts what kind of difference can matter.

\subsection*{9. Choice Can Be Won at Many Stages of the Journey}

The customer does not experience a product as a specification. They experience a sequence: trigger $\to$ discovery $\to$ recognition $\to$ comprehension $\to$ comparison $\to$ approval $\to$ payment $\to$ onboarding $\to$ use $\to$ outcome $\to$ continuation.

\begin{table*}[t]
\centering
\small
\begin{tabularx}{\textwidth}{>{\raggedright\arraybackslash}p{2.8cm} >{\raggedright\arraybackslash}p{5.2cm} >{\raggedright\arraybackslash}X}
\toprule
\textbf{Stage} & \textbf{Typical Failure Point} & \textbf{Relevant Source of Differentiation} \\
\midrule
Discovery & Buyers never encounter your offering & Availability at the exact moment of acute need \\
Recognition & Buyers cannot identify or recall you & Distinctive visual, verbal, and conceptual brand assets \\
Comprehension & Buyers do not understand what you do & Appropriate market category and unambiguous explanation \\
Comparison & Buyers see no relevant superiority & Measurably better outcome, lower risk, or lower friction \\
Approval & The organization vetoes adoption & Clear procurement, verifiable proof, implementation assurance \\
Onboarding & Buyer never reaches first value & Guided setup, compatibility, and structured early support \\
Use & Delivery is poor or unreliable & Operating design, reliability, and service quality \\
Outcome & Customer cannot observe improvement & Rigorous assessment, diagnostic baseline, demonstrated gain \\
Continuation & Buyer regrets purchase or departs & Ongoing capability value, trust, fair terms, and referrals \\
\bottomrule
\end{tabularx}
\caption{Sources of differentiation mapped across the customer decision and usage journey.}
\label{tab:ch8-journey-stages}
\end{table*}

The founder who diagnoses the wrong stage applies the wrong cure. An institution with strong teaching and no awareness needs distribution, not an additional lesson. A software seller losing on security review needs assurance, not more features. A company acquiring customers who churn needs delivery improvement, not another acquisition campaign.

\subsection*{10. The First Synthesis: Differentiation Is an Allocation Decision}

Resources are scarce. A founder cannot produce every possible benefit at world-class levels. Even if modern AI tools make code, marketing, and content cheaper, capital, organizational attention, customer time, authentic expertise, and credible evidence remain tightly constrained.

A strategy of differentiation therefore states not merely \textbf{what we do better}, but \textbf{what we are willing to be worse at or decline entirely} so that the chosen buyer receives an unusually coherent outcome.

This prepares us for positioning, which is the art of making that system of choices meaningful and legible to the person evaluating it.
